% Options for packages loaded elsewhere
\PassOptionsToPackage{unicode}{hyperref}
\PassOptionsToPackage{hyphens}{url}
\documentclass[
]{article}
\usepackage{xcolor}
\usepackage[margin=1in]{geometry}
\usepackage{amsmath,amssymb}
\setcounter{secnumdepth}{-\maxdimen} % remove section numbering
\usepackage{iftex}
\ifPDFTeX
  \usepackage[T1]{fontenc}
  \usepackage[utf8]{inputenc}
  \usepackage{textcomp} % provide euro and other symbols
\else % if luatex or xetex
  \usepackage{unicode-math} % this also loads fontspec
  \defaultfontfeatures{Scale=MatchLowercase}
  \defaultfontfeatures[\rmfamily]{Ligatures=TeX,Scale=1}
\fi
\usepackage{lmodern}
\ifPDFTeX\else
  % xetex/luatex font selection
\fi
% Use upquote if available, for straight quotes in verbatim environments
\IfFileExists{upquote.sty}{\usepackage{upquote}}{}
\IfFileExists{microtype.sty}{% use microtype if available
  \usepackage[]{microtype}
  \UseMicrotypeSet[protrusion]{basicmath} % disable protrusion for tt fonts
}{}
\makeatletter
\@ifundefined{KOMAClassName}{% if non-KOMA class
  \IfFileExists{parskip.sty}{%
    \usepackage{parskip}
  }{% else
    \setlength{\parindent}{0pt}
    \setlength{\parskip}{6pt plus 2pt minus 1pt}}
}{% if KOMA class
  \KOMAoptions{parskip=half}}
\makeatother
\usepackage{color}
\usepackage{fancyvrb}
\newcommand{\VerbBar}{|}
\newcommand{\VERB}{\Verb[commandchars=\\\{\}]}
\DefineVerbatimEnvironment{Highlighting}{Verbatim}{commandchars=\\\{\}}
% Add ',fontsize=\small' for more characters per line
\usepackage{framed}
\definecolor{shadecolor}{RGB}{248,248,248}
\newenvironment{Shaded}{\begin{snugshade}}{\end{snugshade}}
\newcommand{\AlertTok}[1]{\textcolor[rgb]{0.94,0.16,0.16}{#1}}
\newcommand{\AnnotationTok}[1]{\textcolor[rgb]{0.56,0.35,0.01}{\textbf{\textit{#1}}}}
\newcommand{\AttributeTok}[1]{\textcolor[rgb]{0.13,0.29,0.53}{#1}}
\newcommand{\BaseNTok}[1]{\textcolor[rgb]{0.00,0.00,0.81}{#1}}
\newcommand{\BuiltInTok}[1]{#1}
\newcommand{\CharTok}[1]{\textcolor[rgb]{0.31,0.60,0.02}{#1}}
\newcommand{\CommentTok}[1]{\textcolor[rgb]{0.56,0.35,0.01}{\textit{#1}}}
\newcommand{\CommentVarTok}[1]{\textcolor[rgb]{0.56,0.35,0.01}{\textbf{\textit{#1}}}}
\newcommand{\ConstantTok}[1]{\textcolor[rgb]{0.56,0.35,0.01}{#1}}
\newcommand{\ControlFlowTok}[1]{\textcolor[rgb]{0.13,0.29,0.53}{\textbf{#1}}}
\newcommand{\DataTypeTok}[1]{\textcolor[rgb]{0.13,0.29,0.53}{#1}}
\newcommand{\DecValTok}[1]{\textcolor[rgb]{0.00,0.00,0.81}{#1}}
\newcommand{\DocumentationTok}[1]{\textcolor[rgb]{0.56,0.35,0.01}{\textbf{\textit{#1}}}}
\newcommand{\ErrorTok}[1]{\textcolor[rgb]{0.64,0.00,0.00}{\textbf{#1}}}
\newcommand{\ExtensionTok}[1]{#1}
\newcommand{\FloatTok}[1]{\textcolor[rgb]{0.00,0.00,0.81}{#1}}
\newcommand{\FunctionTok}[1]{\textcolor[rgb]{0.13,0.29,0.53}{\textbf{#1}}}
\newcommand{\ImportTok}[1]{#1}
\newcommand{\InformationTok}[1]{\textcolor[rgb]{0.56,0.35,0.01}{\textbf{\textit{#1}}}}
\newcommand{\KeywordTok}[1]{\textcolor[rgb]{0.13,0.29,0.53}{\textbf{#1}}}
\newcommand{\NormalTok}[1]{#1}
\newcommand{\OperatorTok}[1]{\textcolor[rgb]{0.81,0.36,0.00}{\textbf{#1}}}
\newcommand{\OtherTok}[1]{\textcolor[rgb]{0.56,0.35,0.01}{#1}}
\newcommand{\PreprocessorTok}[1]{\textcolor[rgb]{0.56,0.35,0.01}{\textit{#1}}}
\newcommand{\RegionMarkerTok}[1]{#1}
\newcommand{\SpecialCharTok}[1]{\textcolor[rgb]{0.81,0.36,0.00}{\textbf{#1}}}
\newcommand{\SpecialStringTok}[1]{\textcolor[rgb]{0.31,0.60,0.02}{#1}}
\newcommand{\StringTok}[1]{\textcolor[rgb]{0.31,0.60,0.02}{#1}}
\newcommand{\VariableTok}[1]{\textcolor[rgb]{0.00,0.00,0.00}{#1}}
\newcommand{\VerbatimStringTok}[1]{\textcolor[rgb]{0.31,0.60,0.02}{#1}}
\newcommand{\WarningTok}[1]{\textcolor[rgb]{0.56,0.35,0.01}{\textbf{\textit{#1}}}}
\usepackage{graphicx}
\makeatletter
\newsavebox\pandoc@box
\newcommand*\pandocbounded[1]{% scales image to fit in text height/width
  \sbox\pandoc@box{#1}%
  \Gscale@div\@tempa{\textheight}{\dimexpr\ht\pandoc@box+\dp\pandoc@box\relax}%
  \Gscale@div\@tempb{\linewidth}{\wd\pandoc@box}%
  \ifdim\@tempb\p@<\@tempa\p@\let\@tempa\@tempb\fi% select the smaller of both
  \ifdim\@tempa\p@<\p@\scalebox{\@tempa}{\usebox\pandoc@box}%
  \else\usebox{\pandoc@box}%
  \fi%
}
% Set default figure placement to htbp
\def\fps@figure{htbp}
\makeatother
\setlength{\emergencystretch}{3em} % prevent overfull lines
\providecommand{\tightlist}{%
  \setlength{\itemsep}{0pt}\setlength{\parskip}{0pt}}
\usepackage{bookmark}
\IfFileExists{xurl.sty}{\usepackage{xurl}}{} % add URL line breaks if available
\urlstyle{same}
\hypersetup{
  pdftitle={Non-parametric tests},
  hidelinks,
  pdfcreator={LaTeX via pandoc}}

\title{Non-parametric tests}
\author{}
\date{\vspace{-2.5em}}

\begin{document}
\maketitle

\textbf{Inference labs} use the five-step template: Hypothesis →
Visualise → Assumptions → Conduct → Conclude.

\subsection{Learning objectives}\label{learning-objectives}

\begin{itemize}
\tightlist
\item
  Choose the right rank-based test for a given design --- Wilcoxon
  signed-rank, Mann-Whitney \emph{U}, Kruskal-Wallis, sign test.
\item
  Run each in base R and interpret the output.
\item
  State what null hypothesis each test is really addressing (not always
  ``medians are equal'').
\end{itemize}

\subsection{Prerequisites}\label{prerequisites}

Labs 3.4, 4.1.

\subsection{Background}\label{background}

Non-parametric tests make weaker distributional assumptions than the
parametric tests that dominate the preceding labs. They tend to test
\emph{exchangeability} of distributions rather than equality of means,
and they operate on ranks of the data rather than on the raw values.
When the data are heavily skewed, ordinal, or small-sample, the
non-parametric alternatives are often the defensible default.

The workhorse set: Wilcoxon signed-rank for paired continuous data;
Mann-Whitney \emph{U} (the \texttt{wilcox.test} with
\texttt{paired\ =\ FALSE}) for two independent groups; Kruskal-Wallis
for three or more groups; sign test as a cruder paired-data option based
on the signs of differences alone.

A common mistake is to read these tests as ``comparing medians''. They
compare \emph{distributions}. Under additional assumptions --- notably
that the two groups' distributions differ only by a location shift ---
the Mann-Whitney \emph{U} test does estimate a location shift, but that
is a stronger assumption than the test itself requires.

\subsection{Setup}\label{setup}

\begin{Shaded}
\begin{Highlighting}[]
\FunctionTok{library}\NormalTok{(tidyverse)}
\FunctionTok{set.seed}\NormalTok{(}\DecValTok{42}\NormalTok{)}
\FunctionTok{theme\_set}\NormalTok{(}\FunctionTok{theme\_minimal}\NormalTok{(}\AttributeTok{base\_size =} \DecValTok{12}\NormalTok{))}
\end{Highlighting}
\end{Shaded}

\subsection{1. Hypothesis}\label{hypothesis}

Three scenarios:

A. Paired: pre/post in 25 patients, skewed outcome. H0: no shift. B.
Independent two-sample: outcome in two arms, skewed. H0: same
distribution. C. Three-group comparison: outcome across three centres.
H0: same distribution across all three.

\subsection{2. Visualise}\label{visualise}

\begin{Shaded}
\begin{Highlighting}[]
\NormalTok{n\_a }\OtherTok{\textless{}{-}} \DecValTok{25}
\NormalTok{pre }\OtherTok{\textless{}{-}} \FunctionTok{exp}\NormalTok{(}\FunctionTok{rnorm}\NormalTok{(n\_a, }\DecValTok{3}\NormalTok{, }\FloatTok{0.3}\NormalTok{))}
\NormalTok{post }\OtherTok{\textless{}{-}}\NormalTok{ pre }\SpecialCharTok{*} \FunctionTok{exp}\NormalTok{(}\FunctionTok{rnorm}\NormalTok{(n\_a, }\SpecialCharTok{{-}}\FloatTok{0.2}\NormalTok{, }\FloatTok{0.3}\NormalTok{))}
\NormalTok{df\_a }\OtherTok{\textless{}{-}} \FunctionTok{tibble}\NormalTok{(}\AttributeTok{id =} \FunctionTok{seq\_len}\NormalTok{(n\_a), pre, post,}
               \AttributeTok{delta =}\NormalTok{ post }\SpecialCharTok{{-}}\NormalTok{ pre)}

\CommentTok{\# Scenario B: independent two groups, lognormal}
\NormalTok{n\_b }\OtherTok{\textless{}{-}} \DecValTok{30}
\NormalTok{grp }\OtherTok{\textless{}{-}} \FunctionTok{rep}\NormalTok{(}\FunctionTok{c}\NormalTok{(}\StringTok{"A"}\NormalTok{, }\StringTok{"B"}\NormalTok{), }\AttributeTok{each =}\NormalTok{ n\_b)}
\NormalTok{y\_b }\OtherTok{\textless{}{-}} \FunctionTok{c}\NormalTok{(}\FunctionTok{exp}\NormalTok{(}\FunctionTok{rnorm}\NormalTok{(n\_b, }\DecValTok{3}\NormalTok{, }\FloatTok{0.5}\NormalTok{)),}
         \FunctionTok{exp}\NormalTok{(}\FunctionTok{rnorm}\NormalTok{(n\_b, }\FloatTok{3.4}\NormalTok{, }\FloatTok{0.5}\NormalTok{)))}
\NormalTok{df\_b }\OtherTok{\textless{}{-}} \FunctionTok{tibble}\NormalTok{(grp, }\AttributeTok{y =}\NormalTok{ y\_b)}

\CommentTok{\# Scenario C: three groups}
\NormalTok{n\_c }\OtherTok{\textless{}{-}} \DecValTok{25}
\NormalTok{centre }\OtherTok{\textless{}{-}} \FunctionTok{rep}\NormalTok{(}\FunctionTok{c}\NormalTok{(}\StringTok{"X"}\NormalTok{, }\StringTok{"Y"}\NormalTok{, }\StringTok{"Z"}\NormalTok{), }\AttributeTok{each =}\NormalTok{ n\_c)}
\NormalTok{y\_c }\OtherTok{\textless{}{-}} \FunctionTok{c}\NormalTok{(}\FunctionTok{exp}\NormalTok{(}\FunctionTok{rnorm}\NormalTok{(n\_c, }\FloatTok{3.0}\NormalTok{, }\FloatTok{0.4}\NormalTok{)),}
         \FunctionTok{exp}\NormalTok{(}\FunctionTok{rnorm}\NormalTok{(n\_c, }\FloatTok{3.2}\NormalTok{, }\FloatTok{0.4}\NormalTok{)),}
         \FunctionTok{exp}\NormalTok{(}\FunctionTok{rnorm}\NormalTok{(n\_c, }\FloatTok{3.5}\NormalTok{, }\FloatTok{0.4}\NormalTok{)))}
\NormalTok{df\_c }\OtherTok{\textless{}{-}} \FunctionTok{tibble}\NormalTok{(centre, }\AttributeTok{y =}\NormalTok{ y\_c)}
\end{Highlighting}
\end{Shaded}

\begin{Shaded}
\begin{Highlighting}[]
\NormalTok{df\_b }\SpecialCharTok{|\textgreater{}}
  \FunctionTok{ggplot}\NormalTok{(}\FunctionTok{aes}\NormalTok{(grp, y, }\AttributeTok{fill =}\NormalTok{ grp)) }\SpecialCharTok{+}
  \FunctionTok{geom\_boxplot}\NormalTok{(}\AttributeTok{alpha =} \FloatTok{0.6}\NormalTok{, }\AttributeTok{colour =} \StringTok{"grey30"}\NormalTok{) }\SpecialCharTok{+}
  \FunctionTok{labs}\NormalTok{(}\AttributeTok{x =} \ConstantTok{NULL}\NormalTok{, }\AttributeTok{y =} \StringTok{"Outcome"}\NormalTok{) }\SpecialCharTok{+}
  \FunctionTok{theme}\NormalTok{(}\AttributeTok{legend.position =} \StringTok{"none"}\NormalTok{)}
\end{Highlighting}
\end{Shaded}

\subsection{3. Assumptions}\label{assumptions}

Wilcoxon signed-rank: paired observations; differences symmetric around
the null. Mann-Whitney \emph{U}: independent observations within and
between groups; if you want a location-shift interpretation, the two
distributions should have the same shape. Kruskal-Wallis: independent
observations; same-shape assumption extends to all \emph{k} groups.

\subsection{4. Conduct}\label{conduct}

\begin{Shaded}
\begin{Highlighting}[]
\NormalTok{wt\_a }\OtherTok{\textless{}{-}} \FunctionTok{wilcox.test}\NormalTok{(df\_a}\SpecialCharTok{$}\NormalTok{post, df\_a}\SpecialCharTok{$}\NormalTok{pre, }\AttributeTok{paired =} \ConstantTok{TRUE}\NormalTok{,}
                    \AttributeTok{conf.int =} \ConstantTok{TRUE}\NormalTok{)}
\CommentTok{\# Sign test: a simple binomial test on positive differences}
\NormalTok{signs }\OtherTok{\textless{}{-}} \FunctionTok{sum}\NormalTok{(df\_a}\SpecialCharTok{$}\NormalTok{delta }\SpecialCharTok{\textgreater{}} \DecValTok{0}\NormalTok{)}
\NormalTok{sign\_test }\OtherTok{\textless{}{-}} \FunctionTok{binom.test}\NormalTok{(signs, n\_a, }\AttributeTok{p =} \FloatTok{0.5}\NormalTok{)}

\CommentTok{\# B: Mann{-}Whitney U}
\NormalTok{wt\_b }\OtherTok{\textless{}{-}} \FunctionTok{wilcox.test}\NormalTok{(y }\SpecialCharTok{\textasciitilde{}}\NormalTok{ grp, }\AttributeTok{data =}\NormalTok{ df\_b, }\AttributeTok{conf.int =} \ConstantTok{TRUE}\NormalTok{)}

\CommentTok{\# C: Kruskal{-}Wallis}
\NormalTok{kw\_c }\OtherTok{\textless{}{-}} \FunctionTok{kruskal.test}\NormalTok{(y }\SpecialCharTok{\textasciitilde{}}\NormalTok{ centre, }\AttributeTok{data =}\NormalTok{ df\_c)}

\FunctionTok{list}\NormalTok{(}\AttributeTok{wilcox\_paired =}\NormalTok{ wt\_a,}
     \AttributeTok{sign\_test =}\NormalTok{ sign\_test,}
     \AttributeTok{mann\_whitney =}\NormalTok{ wt\_b,}
     \AttributeTok{kruskal\_wallis =}\NormalTok{ kw\_c)}
\end{Highlighting}
\end{Shaded}

Follow-up pairwise comparisons with Holm correction for C:

\begin{Shaded}
\begin{Highlighting}[]

\end{Highlighting}
\end{Shaded}

\subsection{5. Concluding statement}\label{concluding-statement}

\begin{quote}
\textbf{Paired comparison (A).} 25 patients, median change =
\texttt{round(median(df\_a\$delta),\ 2)}; Wilcoxon signed-rank \emph{p}
= \texttt{signif(wt\_a\$p.value,\ 2)}, 95\% CI on the Hodges-Lehmann
estimate \texttt{round(wt\_a\$conf.int{[}1{]},\ 2)} to
\texttt{round(wt\_a\$conf.int{[}2{]},\ 2)}. Sign test \emph{p} =
\texttt{signif(sign\_test\$p.value,\ 2)}.

\textbf{Two-group (B).} Mann-Whitney \emph{p} =
\texttt{signif(wt\_b\$p.value,\ 2)}, 95\% CI
\texttt{round(wt\_b\$conf.int{[}1{]},\ 2)} to
\texttt{round(wt\_b\$conf.int{[}2{]},\ 2)}.

\textbf{Three-group (C).} Kruskal-Wallis χ² =
\texttt{round(kw\_c\$statistic,\ 2)}, df = \texttt{kw\_c\$parameter},
\emph{p} = \texttt{signif(kw\_c\$p.value,\ 2)}. Pairwise Wilcoxon tests
with Holm correction indicated the signal lay in the X vs Z contrast.
\end{quote}

Non-parametric tests are not free lunches. They trade a modest amount of
power (about 5\% vs a \emph{t}-test when the \emph{t}-test's assumptions
hold) for robustness to distributional misspecification.

\subsection{Common pitfalls}\label{common-pitfalls}

\begin{itemize}
\tightlist
\item
  Reading a Wilcoxon as a test of medians. It is a test of
  distributional shift.
\item
  Using a rank test when the data are fine on the raw scale, and
  sacrificing power unnecessarily.
\item
  Forgetting to adjust for multiplicity after a Kruskal-Wallis.
\end{itemize}

\subsection{Further reading}\label{further-reading}

\begin{itemize}
\tightlist
\item
  Hollander M, Wolfe DA. \emph{Nonparametric Statistical Methods}.
\item
  Conover WJ. \emph{Practical Nonparametric Statistics}.
\end{itemize}

\subsection{Session info}\label{session-info}

\begin{Shaded}
\begin{Highlighting}[]

\end{Highlighting}
\end{Shaded}

\subsection{See also --- chapter index}\label{see-also-chapter-index}

\begin{itemize}
\tightlist
\item
  \href{../index.Rmd}{Methods in this chapter}
\item
  \href{../../../ch00-find-your-method.Rmd}{Find your method (Chapter
  0)}
\end{itemize}

\end{document}
